% Proposed replacement for Lemmas pair-five-endpoints and pair-five-reset.
% The seven words are short Hamilton cycles, not endpoint-path tables.

For lists $A,B\subseteq[5]$, write $F(A,B)$ for the graph whose vertices are
$(a,b)\in A\times B$ with $a\ne b$, two vertices being adjacent when one
coordinate is unchanged.  It is useful to regard this as a rook graph with a
(diagonal) matching of cells deleted.  The following is the only finite input
needed by the reset argument.

\begin{lemma}[Edge cycles and common-interface exits]
\label{lem:pair-five-edge-cycles}
Let $|A|,|B|\in\{3,4\}$, after interchanging the lists if necessary.  For
one of the seven types
\[
331,332,333,342,343,443,444,
\]
every edge of $F(A,B)$ lies on a Hamilton cycle.  (For type $333$ the graph
itself is a $6$-cycle.)

Suppose $J=F(A,B)\cap F(A',B')$ is the common state set of one of the
interfaces used in Lemma~\ref{lem:pair-five-reset}: the two lists are
unchanged; one $3$-list is replaced by another $3$-list with one colour
exchanged; or a containing $3$-list and $4$-list are interchanged.  Then
\[
  \forall p\in F(A,B)\quad \exists q\in J\quad pq\in E(F(A,B)).
\]
Consequently, for every $p\in F(A,B)$ there is a $p$--$q$ Hamilton path in
$F(A,B)$ with $q\in J$.
\end{lemma}

\begin{proof}
The seven cycle certificates below use the representatives indicated by the
subscripts.  A word is read cyclically, and $ab$ means $(a,b)$:
\[
\begin{aligned}
 C_{331}&:03\ 04\ 14\ 10\ 20\ 24\ 23\ 13,\\
 C_{332}&:23\ 03\ 01\ 21\ 20\ 10\ 13,\\
 C_{333}&:10\ 20\ 21\ 01\ 02\ 12,\\
 C_{342}&:24\ 04\ 01\ 03\ 13\ 14\ 10\ 20\ 21\ 23,\\
 C_{343}&:12\ 10\ 20\ 21\ 01\ 02\ 03\ 23\ 13,\\
 C_{443}&:32\ 12\ 10\ 20\ 21\ 01\ 02\ 04\ 24\ 14\ 34\ 31\ 30,\\
 C_{444}&:23\ 21\ 20\ 10\ 12\ 02\ 32\ 30\ 31\ 01\ 03\ 13.
\end{aligned}
\]
Each word contains every state of its representative once, and consecutive
(including last--first) states agree in one coordinate.  Apply all colour
permutations preserving the two displayed lists.  For type $332$ also apply
$(a,b)\mapsto(\sigma(b),\sigma(a))$, where $\sigma=(2\ 3)$; this is the
side-exchange symmetry of the representative.  A direct edge-orbit check
shows that the resulting orbit of the displayed cycle has respectively
\[
 14,10,6,20,15,30,24
\]
edges, which are exactly the edge numbers of the seven representative
rook graphs.  Hence every edge lies on a Hamilton cycle.  This is an
edge statement, rather than an endpoint enumeration.

It remains to find a common neighbour in $J$.  If no list changes, use any
row or column neighbour of $p$ (all rows and columns occurring here contain
at least two states).  If, say, the first $3$-list changes to another
$3$-list, put $A_0=A\cap A'$ (so $|A_0|=2$), while $B$ is unchanged and has
size at least $3$.  For $p=(a,b)$, if $a\in A_0$ choose
$d\in B\setminus\{a,b\}$ and put $q=(a,d)$; if $a\notin A_0$, choose
$c\in A_0\setminus\{b\}$ and put $q=(c,b)$.  The choices are nonempty by
$|B|\ge3$ and $|A_0|=2$, respectively.  The case in which the second list
changes is symmetric.  Finally, when a $3$-list $C$ and a containing
$4$-list are interchanged, a state on the $F(C,D)$ side has first coordinate
in $C$, and a state on the $F(D,C)$ side has second coordinate in $C$;
choosing the other coordinate in $C$ gives a neighbour in
$J=F(C,C)$.  Thus every $p$ has a neighbour $q\in J$.  The edge $pq$ lies
on a Hamilton cycle, and deleting that edge gives the required Hamilton
path.
\end{proof}

% Optional structural replacement for the old Hamilton-connected endpoint table.
% It proves the four HC types by row/column clique splicing; only four exceptional
% endpoint words remain, instead of the 39 words for these types.

For a cell $x=(r,c)$ call $r$ its row and $c$ its column.  A row-block path
from $(r,c)$ to $(r,d)$ is obtained by listing all allowed cells in row $r$
with those prescribed first and last cells.  Consecutive row blocks are joined
through a common column.  The analogous column-block construction is obtained
by interchanging the coordinates.

\begin{lemma}[Rook-splice lemma]
\label{lem:rook-splice-hc}
Let $D$ be a matching of forbidden cells, with at most one forbidden cell in
each row and column.  For the four parameter triples
\[
 (|A|,|B|,|D|)=(3,4,2),(3,4,3),(4,4,3),(4,4,4),
\]
the graph $F(A,B)=(A\times B\setminus D)$ with row/column adjacency is
Hamilton-connected.
\end{lemma}

\begin{proof}
We record the elementary bridge argument because it is reusable.  Put
$\epsilon(r)=\{r\}$ when $(r,r)\in D$ and $\epsilon(r)=\varnothing$
otherwise.  If $r_1,\ldots,r_m$ are the rows in the order in which their
blocks are to be traversed, a bridge column between $r_i$ and $r_{i+1}$ may
be chosen from
\[
 S_i=B\setminus(\epsilon(r_i)\cup\epsilon(r_{i+1})).                 \tag{*}
\]
Every $S_i$ has at least two elements.  Choose $c_i\in S_i$ with
$c_i\ne c_{i+1}$; at the first and last bridge also require
$c_1\ne c(x)$ and $c_{m-1}\ne c(y)$.  Listing each row clique from its
incoming bridge to its outgoing bridge (and putting $x$ first and $y$ last)
then gives a spanning $x$--$y$ path.  The same statement holds with rows and
columns interchanged.

For completeness, the small list-choice check is as follows.  If $m=4$, try
the two orders of the middle rows.  For a fixed order, choose $c_2$ first,
then choose $c_1,c_3$ from their two-element lists in (*) avoiding $c_2$ and
the endpoint colour.  A failure forces both endpoint lists to be singletons
and $S_2$ to be their two-element union; swapping the middle rows changes
the two deleted colours in $S_2$ and resolves this obstruction.  If $m=3$,
try the row-block construction and its
column-block transpose.  Failure of both constructions requires the two
endpoint restrictions to reduce the
first and last lists to the same singleton.  With four columns this can only
happen when the rows have respectively three common colours (type $343$), or
two common colours and one outside colour (type $342$).  Interchanging rows
and columns and applying the colour symmetries leaves, respectively, the
three endpoint pairs
\[
 (12,13),\qquad(01,10),\qquad(01,02)                       \tag{343}
\]
and the one pair
\[
 (13,14) .                                                   \tag{342}
\]
All other endpoint pairs are supplied by the block path above.  In type
$443$ and $444$ the $m=4$ choice never fails.

The four exceptional pairs have the following paths (read from left to
right):
\[
\begin{aligned}
343:&\quad 12\ 10\ 20\ 23\ 21\ 01\ 02\ 03\ 13,\\
    &\quad 01\ 02\ 12\ 13\ 03\ 23\ 21\ 20\ 10,\\
    &\quad 01\ 03\ 23\ 21\ 20\ 10\ 13\ 12\ 02;\\
342:&\quad 13\ 10\ 20\ 24\ 21\ 23\ 03\ 01\ 04\ 14.
\end{aligned}
\]
Permuting the common colours (and, for $342$, the two outside columns) gives
all failures listed above.  Thus every endpoint pair has a spanning path.
\end{proof}

\begin{remark}
The old $53$ endpoint words prove a strictly stronger assertion (Hamilton
connectedness of four types and an exact endpoint classification of the other
three).  The reset argument never uses that assertion: it only needs one
common-interface exit for each entrance.  Lemma~\ref{lem:pair-five-edge-cycles}
therefore removes the endpoint table without weakening any subsequent
induction.
\end{remark}
